% =====================================================================
%  4.1 Gjennomsnittlig og momentan vekstfart — fasit
% =====================================================================

\fk{Del 1 — Gjennomsnittlig vekstfart}
\fo{1.1} a) $4$ \ b) $5$ \ c) $2$
\fo{1.2} a) $p(1) = -5$ og $p(3) = -9$ \ b) $\frac{-9 - (-5)}{3 - 1} = \frac{-4}{2} = -2$
\fo{1.3} a) $75$ km/h \ b) $70$ km/h \ c) $70$, $80$ og $60$ km/h. Raskest i den andre timen ($t = 1$ til $t = 2$).
\fo{1.4} a) $3$ \ b) $\frac{0 - 3}{4 - 1} = -1$ \ c) $\frac{3 - 0}{3} = 1$ \ d) $0$ \ e) Nei. $f(1) = f(3) = 3$, men $f(2) = 4$.
\fo{1.5} a) $V(3) = 122\,880$, så $\frac{122\,880 - 240\,000}{3} = -39\,040 \approx -39\,000$ kr per år.
b) $V(5) \approx 78\,643$, så $\frac{78\,643 - 122\,880}{2} \approx -22\,100$ kr per år.
c) $20\,\%$ av en lavere verdi er færre kroner.
\fo{1.6}  Feilen er i $\frac{f(2) - f(6)}{6 - 2}$: $\Delta y$ er regnet feil vei. Riktig er $\frac{18 - (-2)}{6 - 2} = 5$.
\fo{1.7} a) Alltid. $\frac{(kb + m) - (ka + m)}{b - a} = \frac{k(b - a)}{b - a} = k$.
b) Noen ganger. Vekstfarten er $\frac{b^2 - a^2}{b - a} = a + b$, som er positiv bare når $a + b > 0$. På $[1, 2]$ er den $3$, men på $[-3, 2]$ er den $-1$.
c) Noen ganger. $[0, 2]$ gir $2$, mellom $0$ og $4$. $[10, 11]$ gir $21$, ikke mellom $100$ og $121$.
\fo{1.8} a) $f(6) = f(2) + (-3) \cdot 4 = -7$.
b) $f(5) = 25 + 5p$ og $f(1) = 1 + p$. Vekstfarten er $\frac{24 + 4p}{4} = 6 + p = 3$, så $p = -3$.
\fo{Sjekk 1} a) $\frac{10 - 0}{2} = 5$ \ b) $\frac{6 - 12}{17 - 14} = -2$ $^\circ$C per time \ c) At $f$ synker i snitt på intervallet.

\fk{Del 2 — Momentan vekstfart}
\fo{2.1} Rad $h = 0,1$: $1,21$; $0,21$; $2,1$. Rad $h = 0,01$: $1,0201$; $0,0201$; $2,01$. Tallene nærmer seg $2$.
\fo{2.2} a) $\frac{4 - 0}{4 - 2} = 2$ \ b) Voksende. Tangenten har positivt stigningstall, og linja går oppover. \ c) $0$. Tangenten er vannrett.
\fo{2.3} a) Momentan, $70$ km/h. \ b) Gjennomsnittlig, $\frac{270}{3} = 90$ km/h. \ c) Gjennomsnittlig, $\frac{10 - 4}{9 - 6} = 2\ {}^\circ$C per time. \ d) Momentan, $3\ {}^\circ$C per time.
\fo{2.4} a) $h = 1$: $7$ \ $h = 0,1$: $3,31$ \ $h = 0,01$: $3,0301$ \ $h = 0,001$: $3,003001$
b) Vekstfarten i $x = 1$ er omtrent $3$.
c) $h = 0,01$: $12,0601$. $h = 0,001$: $12,006001$. Vekstfarten i $x = 2$ er omtrent $12$.
\fo{2.5} a) $\frac{620}{60} \approx 10,3$ m/s
b) $[20, 30]$: $\frac{120}{10} = 12$ m/s. $[30, 40]$: $\frac{130}{10} = 13$ m/s. Anslag: $12,5$ m/s.
c) a) er farten i snitt over hele etappen. b) er farten i ett øyeblikk, og kan være en helt annen enn snittet.
\fo{2.6} Feilen er i siste overgang: $\frac{0,31}{0,1} = 3,1$, ikke $0,031$. Riktig anslag er $3,1$ (enhet ikke oppgitt).
\fo{2.7} a) $\frac{64 - 4}{6} = 10$
b) $h = 0,1$: $10,1$. $h = 0,01$: $10,01$. Anslag: $10$.
c) Begge er $10$. Eksempel: $[-3, 1]$ gir $-2$, og i $x = -1$ gir $h = 0,01$ tallet $-1,99$. Det ser ut til å gjelde alltid, men eksempler er ikke bevis.
\fo{2.8} a) Omtrent $4$.
b) Nei. Tallene faller fra $5$ via $4,1$ til $4,01$ og nærmer seg $4$, ikke $4,5$.
c) $f(x) = x^2$ passer: $h = 1$ gir $\frac{9 - 4}{1} = 5$, $h = 0,1$ gir $4,1$ og $h = 0,01$ gir $4,01$. Flere funksjoner kan passe.
\fo{Sjekk 2} a) $f(0,1) = 1,01$, så $\frac{1,01 - 1}{0,1} = 0,1$ \ b) $\frac{5 - 1}{2 - 0} = 2$ \ c) Det andre punktet kommer nærmere det første, og sekanten dreier seg inn mot tangenten.

\fk{Ekstra utfordring}
\fo{U1} a) $f(1) = 0$ og $f(3) = 10$, så $5$. $f(-1) = 6$ og $f(2) = 3$, så $\frac{3 - 6}{3} = -1$.
b) $\frac{f(b) - f(0)}{b} = \frac{2b^2 - 3b}{b} = 2b - 3 = 5$, så $b = 4$.
\fo{U2} a) $\frac{3}{2}$, $6$ og $24$. b) Hver ny verdi er $4$ ganger den forrige. $2^{x+2} = 4 \cdot 2^x$, så både $f(b) - f(a)$ og vekstfarten blir $4$ ganger så stor når intervallet flyttes $2$ til høyre.
\fo{U3} a) $T(0) = 80$, $T(5) \approx 55,4$, $T(10) \approx 40,9$. Snitt: $-4,9\ {}^\circ$C per minutt og $-2,9\ {}^\circ$C per minutt.
b) $h = 0,01$: $-3,73$. $h = 0,001$: $-3,73$. Omtrent $-3,73\ {}^\circ$C per minutt.
c) Jo kaldere kaffen er, desto mindre er forskjellen til romtemperaturen, og da går avkjølingen saktere.
\fo{U4} a) Venstre tall er $h$, høyre er sekantens stigningstall. b) Omtrent $0,693$.
c) Med $x = 3$ blir tallet omtrent $5,545$, som er $8$ ganger så mye. Det er fordi $2^{3+h} - 2^3 = 8(2^h - 1)$.
\fo{U5} a) $\frac{a x_2^2 + b x_2 + c - (a x_1^2 + b x_1 + c)}{x_2 - x_1} = \frac{a(x_2^2 - x_1^2) + b(x_2 - x_1)}{x_2 - x_1} = a(x_1 + x_2) + b$.
b) Her er $a = 1$, $b = 0$, $x_1 = 3$ og $x_2 = 3 + h$. Da er vekstfarten $1 \cdot (3 + 3 + h) + 0 = 6 + h$.
\fo{U6} a) $f(1) = 5$. Vi trenger $f(b) = 5$, altså $-b^2 + 6b = 5$, så $b = 5$.
b) Uendelig mange. Fra U5 er vekstfarten $-(p + q) + 6$, som er $0$ for alle $p \neq q$ med $p + q = 6$.
c) Grafen er symmetrisk om $x = 3$. Toppunktet er $(3, 9)$.
d) Eksempel: $f(x) = x(x - 4)$. Da er $f(0) = f(4) = 0$, men $f(1) = -3$. Flere svar er riktige.
\fo{U7} a) Tallene nærmer seg først $0,693$, er best rundt $h = 10^{-8}$, og blir så feil: $0,666$ for $h = 10^{-15}$ og $0,0$ for $h = 10^{-16}$.
b) Maskinen lagrer ca. $16$ siffer. For svært liten $h$ blir differansen $2^h - 1$ rundet av.
